% ============================================================
% ch18.tex —— 第 18 章 变化与积累的桥：微积分基本定理
% ============================================================
\chapter{变化与积累的桥：微积分基本定理}

第 17 章每次算面积都要切一万条长方形，每次都要动用平方和、立方和这类"外援公式"——照这个节奏，$\int_0^1 t^{10}\,dt$ 得先推导十次方和公式。这一章告诉你：\textbf{再也不用切了}。全书最壮丽的定理即将登场——它是连接"变化"（导数）与"积累"（积分）的桥，牛顿与莱布尼茨各自独立发现它，微积分因此正式加冕为"一门学问"。

\section{让面积自己动起来}

\begin{kaiwei}
$y = t^2$ 下方的面积从 $0$ 积到 $1$ 是 $\tfrac13$（上一章挣的）。那积到 $x = 0.5$、$x = 2$、$x = 3$ 呢？把右边界 $x$ 当成\textbf{变量}，"面积"自己怎么变？
\end{kaiwei}

定义\textbf{面积函数}：
\[
F(x) := \int_0^x t^2\,dt \qquad (\text{左边界钉死在 } 0\text{，右边界 } x\text{ 滑动}) .
\]
$F$ 是货真价实的函数——给一个右边界，还你一块面积：$F(0) = 0$（零宽条）；$F(1) = \tfrac13$（上一章的成果）；$F(2) = ?$（先别猜，往下读三段就有答案）。

灵魂问题：\textbf{$F$ 自己有没有导数？}——右边界挪动一点点，面积增长多快？\textbf{一个函数的"函数"求导}，这个句式看着眼晕，其实操作是老三样（差商、化简、取落脚点）。

\section{桥的直觉：一条窄边}

把右边界从 $x$ 挪到 $x + h$。新增的面积是贴着右边界的一条\textbf{窄边}——左边是曲线 $y = t^2$ 在 $x$ 处的取值 $f(x) = x^2$，右边是 $x+h$ 处的取值 $(x+h)^2$，宽度 $h$：

\begin{tikzpicture}[scale=2.2]
  \draw[->] (-0.15,0) -- (1.6,0) node[right] {$t$};
  \draw[->] (0,-0.15) -- (0,1.4) node[above] {$y$};
  \draw[thick,domain=0:1.25,smooth] plot (\x,{\x*\x});
  \draw[fill=cheng!25] (0.8,0) rectangle (1.1,{0.8*0.8});
  \draw[fill=shaohong!40] (0.8,{0.8*0.8}) rectangle (1.1,{1.1*1.1});
  \draw[dashed] (0.8,0) -- (0.8,1.25);
  \node at (0.95,-0.22) {\scriptsize 宽 $h$};
  \node at (1.45,0.55) {\scriptsize 窄边 $\approx f(x)h$};
  \node at (0.42,0.85) {\scriptsize $F(x)$};
  \draw[<->] (1.3,{0.8*0.8}) -- (1.3,{1.1*1.1}) node[midway,right] {\scriptsize 曲线的};
\end{tikzpicture}

新增面积夹在两块之间：竖条 $x^2\cdot h$（矮板凳，高取左端）与 $ (x+h)^2 h$（高板凳，高取右端）。即
\[
x^2 h \;\le\; F(x+h) - F(x) \;\le\; (x+h)^2 h .
\]
（$t^2$ 上升，窄边的高度夹在两端点值之间。）三边同除以 $h$：
\[
x^2 \;\le\; \frac{F(x+h)-F(x)}{h} \;\le\; x^2 + 2xh + h^2 .
\]
$h$ 进场比赛：右边的落脚点是 $x^2$（多项式尾巴压零）。左边钉在 $x^2$——\textbf{夹逼定理}（第 14 章）出庭：中间的差商无处可逃，
\[
F'(x) = x^2 .
\]
面积函数的导数，就是被积函数本身。（这个证明比"窄边近似"的口头版干净：不是"约等于"，是\textbf{两侧真夹逼}。）

\begin{dingli}[微积分基本定理（第一形式）]
设 $f$ 连续（图像不断），定义 $F(x) = \displaystyle\int_a^x f(t)\,dt$，则
\[
F'(x) = f(x).
\]
\end{dingli}

一句话：\textbf{求导与积分互为逆操作}——微分是拆，积分是堆；拆完再堆，物归原主。（严格的"连续"要求保证窄边两端的高度差随 $h$ 消失——$t^2$ 光滑，全书函数全达标。）

\section{后门打开：从此告别切条}

定理立刻兑现为计算神器。想算 $\int_a^b f$，两步走：

\textbf{第一步}：找一个导数恰为 $f$ 的函数 $G$（\textbf{原函数}，土名"反导数"）。\textbf{第二步}：总面积 $= G(b) - G(a)$（端点相减）。

为什么是"端点相减"？三行说清。设 $G$ 是任意一个原函数（$G' = f$）。而 $F(x) = \int_a^x f$ 也是原函数（刚证的）。两个原函数导数处处相等，差 $G - F$ 的导数处处为零——\textbf{导数处处为零的函数，函数值不变}（变化率处处为零 $=$ 一动不动；严格版是大学的"中值定理"，直觉版人人可用）。于是 $G(x) - F(x) = C$（常数），即 $G(x) = F(x) + C$。代入端点：
\[
G(b) - G(a) = \big[F(b) + C\big] - \big[F(a) + C\big] = F(b) - F(a) = \int_a^b f .
\]
常数 $C$ 相减时\textbf{自动抵消}——这就是"原函数有无穷多个（差常数），但端点相减时随便挑一个都行"的原因。

\begin{dingli}[微积分基本定理（第二形式 / 牛顿--莱布尼茨公式）]
若 $G' = f$，则
\[
\int_a^b f(t)\,dt = G(b) - G(a) =: G(t)\Big|_a^b .
\]
（$\Big|$ 读作"代入相减"。）
\end{dingli}

\section{实战七连：后门的威力}

原函数哪里找？\textbf{把幂法则倒着抄}。谁求导得 $t^2$？——$\tfrac{t^3}{3}$（求导：$3\cdot\tfrac{t^2}{3} = t^2$ \checkmark）。逐题验算：

\paragraph{一（清旧账）。}$\displaystyle\int_0^1 t^2\,dt = \frac{t^3}{3}\Big|_0^1 = \frac13 - 0 = \frac13$。\textbf{上一章切一万条的心血，两行结束。}

\paragraph{二（清悬案）。}$\displaystyle\int_0^2 t^2\,dt = \frac{t^3}{3}\Big|_0^2 = \frac83$。开胃问题里"$F(2) = ?$"兑现：$\tfrac83$。

\paragraph{三（通杀）。}$\displaystyle\int_0^1 t^n\,dt = \frac{t^{n+1}}{n+1}\bigg|_0^1 = \frac{1}{n+1}$，对一切 $n$ 成立。第 17 章烧脑时刻"平方和、立方和公式接力"的苦工，一个公式全体终结——$n = 100$ 也只是心算 $\tfrac1{101}$。

\paragraph{四（带常数）。}$\displaystyle\int_0^3 (2t + 5)\,dt = \left[t^2 + 5t\right]_0^3 = 9 + 15 = 24$。

\paragraph{五（物理：落体）。}$v = gt$，下落 $T$ 秒的路程：
\[
\int_0^T gt\,dt = \frac{g t^2}{2}\Big|_0^T = \frac12 gT^2 .
\]
\textbf{伽利略的 $s = \tfrac12 g t^2$ 从"堆积瞬间速度"里自动长出来}——他从斜面实验里猜的规律，在这里是三行代数的必然。取 $g = 9.8$、$T = 3$：$44.1$ 米。

\paragraph{六（物理：刹车）。}$v = 10 - t$（米/秒），停表 $T = 10$：
\[
\int_0^{10}(10 - t)\,dt = \left[10t - \frac{t^2}2\right]_0^{10} = 100 - 50 = 50\ \text{米}.
\]
（速度--时间图上这块区域是底 $10$ 高 $10$ 的三角形，$\tfrac12\times10\times10 = 50$ \checkmark——\textbf{积分与几何互相验算}，两套语言一台戏。）

\paragraph{七（分段）。}$\displaystyle\int_0^2 |1 - t|\,dt$：绝对值在 $t = 1$ 处折一下，分段积分——$\int_0^1 (1-t)dt + \int_1^2(t-1)dt = \tfrac12 + \tfrac12 = 1$。（几何：两个底 $1$ 高 $1$ 的三角形。折线不影响后门，只影响"分几段"。）

\begin{cidian}
土名对照：反导数 $=$ \textbf{原函数}（antiderivative）；"端点相减"$=$ \textbf{牛顿--莱布尼茨公式}（Newton--Leibniz formula）；$\displaystyle\int f\,dx = F + C$（不带边界的版本）叫\textbf{不定积分}——"$+C$"就是那个自动抵消但又不能忘写的常数。微积分三件套配齐：\textbf{导数}（变化率）、\textbf{积分}（积累）、\textbf{基本定理}（互逆）。
\end{cidian}

\section{为什么这座桥价值连城}

没有桥的世界里，导数与积分是两门手艺：一个管"瞬时"，一个管"总量"，各自为战，每题重切。有了桥：

\textbf{计算上}：积分难题（切一万条）批量转化为求导的逆运算（查表倒抄）——积分从"重体力劳动"变成"查字典"。整本《积分表》之所以能存在，全靠这座桥。

\textbf{哲学上}：宇宙的两大类测量——\textbf{速率}（某刻多快）与\textbf{总量}（一共多少）——被证明是同一枚硬币的两面。位移与速度、电量与电流、水量与流量、人口与出生率……所有"流量 $-$ 库存"式的物理对，背后都是这座桥。物理学家写下的每一个方程，几乎都在桥上散步（牛顿第二定律 $F = ma$ 是"位移的二阶导数 $=$ 力"，解出轨道要靠积分过桥）。

\begin{lishi}
牛顿（1643--1727）在 1666 年（伦敦大瘟疫休学回乡的"奇迹年"，23 岁）发展"流数术"——他管变量叫"流量"、变化率叫"流数"，积分是"流量的逆问题"。莱布尼茨（1646--1716）独立于 1675 年前后得到同等结果，并发明了今天通用的记号 $\int$、$dx$、$\dfrac{dy}{dx}$——记号的胜利如此彻底，以至于"莱布尼茨记号好不好用"本身就是数学史的重要章节。两人后半生陷入优先权大战，英德两国数学界为此断交百年（英国坚持牛顿的点记号，元气大伤）。今天的裁决是：\textbf{两人独立发明，莱布尼茨的记号更优}。而"基本定理"这个名字提示我们：它不是微积分的一件产品，是微积分的\textbf{地基}。
\end{lishi}

\begin{dongshou}
\begin{itemize}
  \yisi 用后门算：$\displaystyle\int_0^3 t^2\,dt$（答 $9$）；$\displaystyle\int_1^4 \frac{1}{t^2}\,dt$（原函数 $-\tfrac1t$，答 $\tfrac34$）；$\displaystyle\int_0^2 (3t^2 + 2t)\,dt$（答 $12$）。
  \yier 速度 $v = 2t$，物体走完前 $9$ 米用了几秒？（$\int_0^T 2t\,dt = T^2 = 9$，$T = 3$ 秒。注意开方取正根——时间非负。）
  \yier 求导验证你是"倒抄幂法则"而非死记：先写 $\displaystyle\int t^4\,dt = \tfrac{t^5}{5} + C$，再对右边求导（幂法则）确认回到 $t^4$。同法造三条自己的积分公式（$t^{10}$、$\sqrt t$、$\tfrac1{t^2}$ 各一条，指数加一放分母）。
  \yisi 用几何直觉秒答（不动笔）：$\displaystyle\int_0^4 3\,dt = ?$（高 $3$ 宽 $4$ 的矩形，$12$）；再用后门验算（原函数 $3t$）。
\end{itemize}
\end{dongshou}

\begin{shaonao}
\textbf{把桥走反一遍。}第 15 章的机器是"路程求导得速度"；本章的机器是"速度积分得路程"。现在把两者串成一台"函数变换器"：对 $f(t) = t^2$，定义
\[
F(x) = \int_0^x t^2\,dt = \frac{x^3}{3}, \qquad\text{于是}\qquad F'(x) = x^2, \qquad F''(x) = 2x .
\]
一阶导数还原回自己，二阶导数是 $2x$。请回答三问：

(1) $\dfrac{d}{dx}\displaystyle\int_0^x t^{10}\,dt = ?$（不必算积分！桥的答案一行。）

(2) $\displaystyle\int_0^x \mathrm{e}^t\,dt = ?$——等等，$\mathrm{e}$ 还没发明（下一章）。换一个：$\displaystyle\int_0^x \frac{1}{(1+t)^2}\,dt = ?$（提示：$\left(\tfrac{1}{1+t}\right)' = -\tfrac{1}{(1+t)^2}$，链式法则倒抄；注意负号，答 $1 - \tfrac1{1+x}$。）

(3) 深水区预告：什么样的函数 $f$ 恰好是某个函数的导数（从而"可积分回原函数"）？连续的都行（第一形式保证 $F$ 存在）；但\textbf{原函数能不能写成初等公式}是另一回事——$\mathrm{e}^{-t^2}$（概率论的台柱）连续、可积，其原函数却\textbf{证明没有}初等表达式（刘维尔定理）。所以积分表永远收不全：\textbf{积分存在}与\textbf{积分算得出}，是两个世界的事。这道题的答案，将在你未来遇到"数值积分"时正式接站。
\end{shaonao}

\begin{chengshi}
"导数处处为零 $\Rightarrow$ 函数为常数"严格版依赖连续性与中值定理（大学分析第一课的压轴）；我们的直觉版（一动不动）对全书场景（多项式、光滑函数）完全够用。"连续函数的变上限积分存在且可导"的严格证明（一致连续性）同样放行大学。另外请注意桥的方向性：先积分后求导必回原函数（第一形式），先求导后积分差个常数（$\int f' = f + C$）——"互逆"带一个 $C$ 的余数，不定积分的"$+C$"就是它的签名。
\end{chengshi}

拆与堆的闭环完成。篇三还剩一位主角没出场：一个自带来世背景的数——它管着复利、人口与衰变，还是下一章素数普查的标尺。有请 $\mathrm{e}$。
